\section{Descripción de la Configuración}

\subsection{Configuración del Direccionamiento IPv4}

El plan de direccionamiento IPv4 se construyó a partir del bloque asignado
\texttt{172.16.192.0/19}, aplicando segmentación jerárquica mediante VLSM para
ajustar el tamaño de cada subred a la cantidad real de hosts requeridos en cada
área. Con el crecimiento previsto del 10\% sobre la cantidad de usuarios, LAN 1
requiere capacidad para 936 hosts, LAN 2 para 660 hosts y LAN 3 para 221 hosts.
Por esta razón, LAN 1 y LAN 2 se asignaron con prefijo \texttt{/22}, mientras que
LAN 3 se asignó con prefijo \texttt{/24}, lo que garantiza suficiente espacio
para los equipos finales y permite crecimiento adicional dentro de la misma red.

La red de servidores se asignó en una subred independiente de tamaño reducido,
\texttt{172.16.202.0/29}, con el fin de alojar los servicios centrales de la
infraestructura sin desperdiciar direcciones. De igual forma, los enlaces entre
routers se configuraron con subredes \texttt{/30}, dado que solo requieren dos
direcciones utilizables por enlace. Este criterio reduce el consumo del espacio
de direccionamiento y mantiene una organización clara entre redes de usuarios,
red de servicios y enlaces de interconexión.

La tabla~\ref{tab:ipv4} resume la asignación propuesta para toda la topología.

\begin{table}[H]
\centering
\caption{Plan de direccionamiento IPv4 para el Grupo 1}
\label{tab:ipv4}
\begin{tabularx}{\textwidth}{l l l l l l}
\toprule
Segmento & Red & Prefijo & Primer host & Último host & Broadcast \\
\midrule
LAN 1 & 172.16.192.0 & /22 & 172.16.192.1 & 172.16.195.254 & 172.16.195.255 \\
LAN 2 & 172.16.196.0 & /22 & 172.16.196.1 & 172.16.199.254 & 172.16.199.255 \\
LAN 3 & 172.16.200.0 & /24 & 172.16.200.1 & 172.16.200.254 & 172.16.200.255 \\
Servidores & 172.16.202.0 & /29 & 172.16.202.1 & 172.16.202.6 & 172.16.202.7 \\
Enlace R1-R2 & 172.16.202.8 & /30 & 172.16.202.9 & 172.16.202.10 & 172.16.202.11 \\
\bottomrule
\end{tabularx}
\end{table}

\subsection{Configuración del Direccionamiento IPv6}

Para IPv6 se utilizó el prefijo asignado \texttt{2001:dbe:bef9::/48}, del cual se
derivaron subredes \texttt{/64} para cada segmento de la topología. Esta decisión
permite mantener una estructura uniforme en todos los enlaces y simplifica la
configuración de autoconfiguración en los hosts finales mediante SLAAC. A diferencia
de IPv4, en IPv6 no es necesario ajustar el tamaño de subred según la cantidad de
hosts, por lo que se adoptó una asignación homogénea por segmento.

La tabla~\ref{tab:ipv6} presenta los prefijos utilizados para cada red y la
dirección sugerida para la interfaz del router en cada segmento.

Adicionalmente, la tabla~\ref{tab:servicios} describe los servicios de red implementados en la topología, incluyendo los servidores web y de correo, junto con sus direcciones IP, protocolos utilizados y puertos asociados. Esta información permite identificar claramente los recursos disponibles en la red y los mecanismos de comunicación empleados.

Por otra parte, en la tabla~\ref{tab:hosts} se muestra la asignación de direcciones IPv4 e IPv6 en los hosts finales, donde se evidencia la distribución organizada de direcciones dentro de cada subred, así como la configuración de sus respectivas puertas de enlace para garantizar la conectividad entre redes.

\begin{table}[H]
\centering
\caption{Plan de direccionamiento IPv6 para el Grupo 1}
\label{tab:ipv6}
\begin{tabularx}{\textwidth}{l l l}
\toprule
Segmento & Prefijo IPv6 & Gateway sugerido \\
\midrule
LAN 1 & 2001:dbe:bef9:1::/64 & 2001:dbe:bef9:1::1 \\
LAN 2 & 2001:dbe:bef9:2::/64 & 2001:dbe:bef9:2::1 \\
LAN 3 & 2001:dbe:bef9:3::/64 & 2001:dbe:bef9:3::1 \\
Servidores & 2001:dbe:bef9:10::/64 & 2001:dbe:bef9:10::1 \\
Enlace R1-R2 & 2001:dbe:bef9:ff01::/64 & 2001:dbe:bef9:ff01::1 \\
\bottomrule
\end{tabularx}
\end{table}

\begin{table}[H]
\centering
\caption{Servicios de red y puertos utilizados}
\label{tab:servicios}
\begin{tabularx}{\textwidth}{l l l l l}
\toprule
Servicio & IPv4 & IPv6 (manual) & Puerto & Tipo \\
\midrule
Servidor Web & 172.16.202.6 & 2001:dbe:bef9:10::10 & 80 & HTTP \\
Servidor Correo (SMTP) & 172.16.202.5 & 2001:dbe:bef9:10::11 & 25 & SMTP \\
Servidor Correo (POP3) & 172.16.202.5 & 2001:dbe:bef9:10::11 & 110 & POP3 \\
\bottomrule
\end{tabularx}
\end{table}

\begin{table}[H]
\centering
\caption{Asignación de direcciones IP en hosts}
\label{tab:hosts}
\begin{tabularx}{\textwidth}{l l l l l}
\toprule
Dispositivo & Red & IPv4 & IPv6 (manual) & Gateway \\
\midrule
PC1-1 & LAN 1 & 172.16.192.253 & 2001:dbe:bef9:1::10 & 172.16.192.1 \\
PC1-2 & LAN 1 & 172.16.192.254 & 2001:dbe:bef9:1::11 & 172.16.192.1 \\
PC2-1 & LAN 2 & 172.16.196.253 & 2001:dbe:bef9:2::10 & 172.16.196.1 \\
PC2-2 & LAN 2 & 172.16.196.254 & 2001:dbe:bef9:2::11 & 172.16.196.1 \\
PC3-1 & LAN 3 & 172.16.200.253 & 2001:dbe:bef9:3::10 & 172.16.200.1 \\
PC3-2 & LAN 3 & 172.16.200.254 & 2001:dbe:bef9:3::11 & 172.16.200.1 \\
\bottomrule
\end{tabularx}
\end{table}

\subsection{Configuración de los Equipos de Red}

Los routers de la topología se configuran con direccionamiento estático en
ambas familias de protocolos. Cada interfaz conectada a una LAN recibe una
dirección de puerta de enlace dentro del rango correspondiente, mientras que
las interfaces de interconexión entre routers utilizan las subredes punto a
punto definidas anteriormente. Esta organización facilita la administración de
rutas y permite identificar con claridad el papel de cada interfaz dentro de la
topología.

En los dispositivos de acceso y conmutación se habilitan contraseñas locales y
acceso remoto seguro mediante SSHv2. El protocolo Telnet no se utiliza, debido
a que transmite credenciales sin cifrado. Adicionalmente, se limita el acceso
administrativo únicamente a las direcciones internas de la red, de forma que
la configuración remota solo pueda realizarse desde los segmentos autorizados.

\subsection{Configuración de los Servidores y Despliegue Físico}

Aunque el diseño lógico prevé un gran clúster de servidores para atender la demanda total teórica, el despliegue práctico para la validación de esta topología se realiza mediante un esquema distribuido de dos máquinas virtuales (VMs) corriendo Ubuntu 26.04 . Cada máquina virtual fue aprovisionada con 3 núcleos de procesamiento (vCPUs) y 8 GB de memoria RAM, garantizando recursos suficientes para la operación fluida de los servicios en un entorno representativo.

Ambos servidores cuentan con direcciones IP fijas dentro de la subred \texttt{172.16.202.0/29} (con equivalentes IPv6 en \texttt{2001:dbe:bef9:10::/64}):
\begin{itemize}
    \item \textbf{VM 1 (Servidor Web y Portal):} Aloja el portal React/Vite servido mediante Nginx, así como el \textit{backend} en Python (FastAPI) que provee la API REST y los WebSockets.
    \item \textbf{VM 2 (Servidor de Correo y Chat):} Dedicada a los servicios persistentes en segundo plano, corriendo el sistema de mensajería desarrollado en C, además de los demonios Postfix y Dovecot para el correo corporativo.
\end{itemize}

La asignación fija es necesaria para asegurar que los clientes de todas las LAN puedan localizar los servicios de forma consistente en los registros DNS sin depender de mecanismos dinámicos.

El aseguramiento (\textit{hardening}) del SO en estos servidores se orquesta de manera automatizada mediante nuestros scripts de despliegue. En lugar de configuraciones iterativas y manuales, los parámetros de red, el endurecimiento de SSHv2, las contramedidas perimetrales usando UFW y los permisos granulares para la creación de usuarios o encendido de demonios, son inyectados mediante los ejecutables durante la etapa de provisión. La descripción completa de estas políticas de seguridad estructurales se profundiza en las decisiones de diseño.

\subsection{Configuración de los Clientes}

Los equipos de usuario final reciben su configuración de red de manera dinámica.
En IPv4 se emplea DHCP para distribuir dirección IP, máscara, puerta de enlace
y servidor DNS. En IPv6 se utiliza SLAAC, de modo que los clientes generan su
propia dirección a partir del prefijo anunciado por el router en cada segmento.
Esta estrategia disminuye la carga administrativa y facilita la puesta en marcha
de un número elevado de clientes en laboratorio.

El cliente de mensajería instantánea desarrollado en Windows se conecta al
servidor Linux mediante la dirección asignada en la red de servicios, y utiliza
la interfaz de terminal para interactuar con el sistema. El soporte de historial
local, listado de usuarios conectados y salas de chat se valida directamente
sobre la conectividad establecida entre los segmentos de la red.

